\documentclass[12pt]{article}
\usepackage[utf8]{inputenc}
\usepackage[T1]{fontenc}
\usepackage{hyperref}
\title{Wik-Mungkan Language Resources for K-12 Curriculum Development}
\author{Compiled by Claude Sonnet 5.5 (claude-sonnet-5-5) for Andr\'as Kornai}
\date{}
\begin{document}
\maketitle

\section{Introduction}
Wik-Mungkan (ISO 639-3: wim; Glottocode wikm1247) is a Pama-Nyungan language of Cape York, Queensland, Australia, associated with the community of Aurukun. This report is based on a small number of web searches. It lists only resources that appeared in the search results. Little ready-made school material turned up, and none of it is a curriculum.

\section{Primary Educational Resources}
No primers, textbooks or school programmes for Wik-Mungkan appeared in the search results. A K-12 curriculum would therefore have to be built from the reference works below and from work with the Aurukun community.

\section{Linguistic and Reference Materials}
\begin{itemize}
\item \emph{A Grammar of Wik-Mungkan}, Alan Everton Ray's 2021 Monash University PhD thesis. It is the most comprehensive grammar found. It draws on historical oral recordings, earlier analyses and recent conversations with elders in Aurukun. \url{https://bridges.monash.edu/articles/thesis/A_Grammar_of_Wik-Mungkan/15049146}
\item \emph{Dictionary and source book of the Wik-Mungkan language} (Kilham et al., SIL, 1986), recorded in a bibliographic database. \url{https://lexbib.elex.is/wiki/Item:Q28882}
\item Pama Language Centre, ``Wik-Mungkan Linguistic Notes and Sources''. A reading list that includes Kilham on discourse, compounds and word order, Godfrey on verb morphology, and Sayers on stress and pitch. \url{https://www.pamacentre.org.au/wik-mungkan-linguistic-notes-and-sources/}
\item Wikipedia overview of the language. \url{https://en.wikipedia.org/wiki/Wik_Mungkan_language}
\item Glottolog reference entry for Ray (2021). \url{https://glottolog.org/resource/reference/id/589191}
\end{itemize}

\section{Religious and Community Resources}
The Pama Language Centre appears to be the main community-facing organisation. It states that NAATI-accredited Wik-Mungkan interpreters can be contacted through it. \url{https://www.pamacentre.org.au/wik-mungkan-2/}

\section{Digital and Audio Resources}
A search summary said that the dictionary has been published online by AuSIL as the Wik Mungkan-English Interactive Dictionary. No URL for it appeared in the results, so it is not linked here. The Pama Centre notes that open-access archival resources are available through AIATSIS. No specific archive URL was seen.

\section{Curriculum Development Context}
Wik-Mungkan is spoken by all adults in the community but not by all young people, so there is still a living speaker base. It is not known to be taught in schools. A curriculum would need to be developed together with Aurukun elders and the Pama Language Centre. The Ray grammar and the Kilham dictionary are the starting lexical and grammatical sources. Audio from archives would need to be located separately.

\end{document}